\section{Planificación y elementos curriculares}

\subsection{Objetivos didácticos}
Al finalizar, el alumnado será capaz de:
\begin{enumerate}[label=\textbf{O\arabic*.},leftmargin=2.4em]
  \item Representar vectores y operar con ellos en GeoGebra, interpretando geométricamente cada operación (CE7).
  \item Formular conjeturas sobre las razones trigonométricas y demostrarlas (CE3).
  \item Resolver triángulos con los teoremas del seno y del coseno y comprobar la validez de las soluciones (CE1, CE2).
  \item Relacionar vectores, trigonometría y geometría analítica en las ecuaciones de la recta, los ángulos y las distancias (CE5).
  \item Modelizar un problema real de medida y comunicar el método con rigor (CE1, CE7, CE8).
  \item Conectar el modelo con la energía solar y justificar una propuesta para el centro (CE2, CE6, CE9).
\end{enumerate}

\subsection{Competencias, criterios y saberes}
Se trabajan ocho de las nueve competencias específicas de Matemáticas~I (todas salvo la CE4), con sus criterios. De los descriptores del perfil de salida que el decreto conecta con cada competencia \parencite[pp.~321--324]{decreto64} se priorizan solo los que la secuencia trabaja con una actuación y una evidencia concretas. Los saberes son contenidos literales de los bloques A, B, C, D y F \parencite[pp.~325--328]{decreto64}. Las relaciones completas están en la sección~\ref{sec:matrices}.

\subsection{Metodología}
La secuencia combina tres modos de trabajo, cada uno con una función distinta:
\begin{itemize}
  \item \textbf{Sesiones magistrales (S3, S4 y S6--S11).} Cada bloque de contenidos se reparte en dos sesiones, con exposiciones de 20 a 30 minutos, para que demostrar y practicar quepa en el tiempo real de clase. El profesor expone de forma ordenada los contenidos, demuestra los resultados y resuelve ejemplos; después el alumnado practica con ejercicios y cierra con una \emph{pieza del reto}: diez minutos en equipo para aplicar lo aprendido al propio centro. Así el contexto está presente desde las primeras sesiones y no solo al final. El decreto recuerda que la ejercitación sigue siendo fundamental para aprender conceptos y métodos básicos antes de afrontar la resolución de problemas \parencite[p.~321]{decreto64}.
  \item \textbf{Talleres de GeoGebra (S2, S5 y S12).} El alumnado construye, manipula y conjetura. En vectores y trigonometría, el taller precede a la exposición, de modo que la sesión magistral formaliza lo observado. La herramienta actúa como colaboradora que ayuda a comprobar cálculos y conjeturas, no como caja negra que sustituye al cálculo \parencite[p.~321]{decreto64}.
  \item \textbf{Aprendizaje basado en un reto, con trabajo cooperativo (S1 y S14--S20).} Equipos con cuatro roles rotatorios (coordinación y Kanban, medición, modelo, comunicación)\soloProducto{, descritos en la sección~\ref{sec:cooperativo},} resuelven una tarea real. Esto responde al principio de fomentar la capacidad de aprender por sí mismos, de trabajar en equipo y de aplicar métodos de investigación apropiados \parencite[art.~4.1, p.~16]{decreto64}.
\end{itemize}
Los verbos de los criterios operativizados siguen la taxonomía revisada de Bloom \parencite{anderson2001} y avanzan desde la aplicación de procedimientos hasta la creación de un modelo propio.

\subsection{Temporalización}
Veinte sesiones de 55 minutos en cinco semanas, entre la segunda quincena de febrero y la tercera semana de marzo. La duración mínima de la guía es de tres semanas; aquí se amplía para que quepan los contenidos del bloque, el trabajo de campo y una defensa por equipo con tiempo suficiente.

\begin{tablaSAlarga}{colspec={Q[c,wd=26pt]Q[l,wd=52pt]Q[l,wd=84pt]X[l]Q[l,wd=44pt]},
  row{1}={bg=urjc, fg=white, font=\footnotesize\bfseries}}
  Ses. & Fase & Tipo & Contenido & Objetivos \\
  1 & Inicio & Reto & Reto, cuestionario diagnóstico, equipos y presentación del socio eTwinning & -- \\
  2 & Desarrollo & Taller (actividad) & Kanban; explorar vectores en GeoGebra & O1 \\
  3 & Desarrollo & \textbf{Magistral} + ejercicios & Vectores: estructura, módulo y coordenadas & O1 \\
  4 & Desarrollo & \textbf{Magistral} + ejercicios & Producto escalar, ángulo y proyección & O1 \\
  5 & Desarrollo & Taller (actividad) & Circunferencia goniométrica: conjeturas & O2 \\
  6 & Desarrollo & \textbf{Magistral} + ejercicios & Razones de un ángulo cualquiera, identidades, ángulo suma y diferencia & O2 \\
  7 & Desarrollo & \textbf{Magistral} + ejercicios & Ángulo doble y mitad; ecuaciones trigonométricas & O2 \\
  8 & Desarrollo & \textbf{Magistral} + ejercicios & Teoremas del seno y del coseno & O3 \\
  9 & Desarrollo & \textbf{Magistral} + ficha & Caso ambiguo y altura inaccesible & O3 \\
  10 & Desarrollo & \textbf{Magistral} + ejercicios & Ecuaciones de la recta y posiciones relativas & O4 \\
  11 & Desarrollo & \textbf{Magistral} + ejercicios & Ángulos entre rectas y distancias & O4 \\
  12 & Desarrollo & Taller (actividad) & Rayo y sombra en GeoGebra (F4) & O4, O5 \\
  13 & Desarrollo & Taller (actividad) & Refuerzo y ampliación por criterio; mapa conceptual; aprobación del método & O1--O5 \\
  14 & Desarrollo & Tarea & Medición en el patio al mediodía solar; datos eTwinning & O5, O6 \\
  15 & Desarrollo & Tarea & Modelo de sombras y análisis del error & O5, O6 \\
  16 & Desarrollo & Tarea & Propuesta de ampliación, comparación con el socio y coevaluación & O5, O6 \\
  17 & Desarrollo & Tarea & Mejora del libro y ensayo de la defensa & O5, O6 \\
  18 & Cierre & Evaluación & Examen individual y retroalimentación & O1--O4 \\
  19 & Cierre & Defensa & Defensa ante la dirección (cuatro equipos) y coevaluación & O5, O6 \\
  20 & Cierre & Defensa & Defensa (tres equipos), coevaluación, metacognición y cierre eTwinning & O5, O6 \\
\end{tablaSAlarga}
